Un fil de cuivre de \qty{\v{L}}{m} de longueur a un diamètre de \qty{\v{d}}{mm}. La résistivité du cuivre est de \qty{1.68e-8}{\ohm.m}. On applique une différence de potentiel de \qty{\v{V}}{V} entre ses extrémités.

\begin{enumerate}[label=\alph*)]
    \item Déterminez la résistance du fil.
    \item Déterminez le courant qui circule dans le fil.
    \item Déterminez la densité de courant dans le fil et la puissance dissipée.
\end{enumerate}
